\scanpage{72}
\begin{flushright}\handwritten{22.7.961}\end{flushright}

\begin{center}Comments on the article by A.Delzant on quadratic forms.\end{center}

From the point of view of exposition;

a) It is very confusing to pretend to treat at the same time characteristic $\handwritten{\ne}2$ and characteristic 2. For the moment, even the structures that come into play are of different kinds, and the treatment of the two cases is really different. So if you do not arrive at a satisfactory common theory, drop characteristic 2, or pu\handwritten{t} in a separate section what\deleted{\unsure{x}} you can say about it.

b) Contrary to what a reading of the article suggests, the role of the $\handwritten{\lambda}$-structures in the theory is practically nil, probably due to the fact that the natural filtration of the ring of classes of quadratic forms is none other than that defined by the powers of the augmentation ideal. One trivially recovers the $\handwritten{\lambda}$-structure from the explicit form of the said ring, and it suffices to do so at the end.

\deleted{\unsure{a)|c)}} From the mathematical point of view, I suggest as a reasonable minimum programme the following, which is within immediate reach:

1) \uline{Structure of the ring} $Q(A)$ \uline{of classes of quadratic forms over} $A$, where for the moment $A$ is going to be a field of characteristic $\handwritten{\ne}2$. For $a\handwritten{\in}A^{\handwritten{*}}$, let $\handwritten{\varphi}_a$ be the class of the quadratic form $\handwritten{at^2}_{a}$, it depends only on $a\bmod A^{\handwritten{*2}}$, hence a homomorphism from
\[
D(A)=A^{\handwritten{*}}/A^{\handwritten{*2}}=H^{\handwritten{1}}(A,\underline{Z}/2\underline{Z})
\]
into $Q^{\deleted{\unsure{2|1}}}(A)^{\handwritten{*}}$, also denoted $a\handwritten{\to\varphi}_a$, hence a homomorphism
\[
\handwritten{\varphi:}\ \underline{Z}(D)\handwritten{\longrightarrow}Q(A)
\tag*{\handwritten{(*)}}
\]
where $\underline{Z}(D)$ is the group algebra of $D=D(A)$, with coefficients in $\underline{Z}$. One will then prove the following theorem:

\uline{The preceding homomorphism is surjective}, \uline{its kernel is} \deleted{formed of the} \uline{the subgroup generated by the} $e_a+e_b-e_{a'}-e_{b'}$, \uline{for} $a,b,a',b'\handwritten{\in}D$ \uline{such that} $\handwritten{\varphi}_a+\handwritten{\varphi}_b=\handwritten{\varphi}_{a'}+\handwritten{\varphi}_{b'}$. According to a theorem of Witt, this last relation is also equivalent to saying that the forms on the two sides have the same invariants, which is written more simply
\reversemarginpar
\margin{\handwritten{\begin{tikzpicture}[baseline=0pt]\draw[gray] (0.4,-1.7)--(0.37,0.2)--(0.4,1.6);\node at (0,0) {?};\end{tikzpicture}}}
\normalmarginpar
\[
a+b=a'+b'\ ,\quad a.b=a'.b'
\]
where $D$ is written additively, and where the product $a.b$ is the cup product, with values in
\[
B=H^{\handwritten{2}}(A,\underline{Z}/2\underline{Z})
\]
% archive sheet number: 71
